% Aula 05 --- poda por custo-complexidade e o papel de alpha.
Construir uma árvore tem duas etapas, e elas usam objetivos diferentes de
propósito. O crescimento escolhe cortes que mais reduzem o RSS; a poda escolhe uma
subárvore $T\subseteq T_0$ minimizando
\[
  C_\alpha(T) \;=\; \sum_{k=1}^{\abs{T}}\ \sum_{i:\;\X_i\in R_k}
    \big(Y_i-\widehat y_{R_k}\big)^2 \;+\; \alpha\,\abs{T},
\]
onde $\abs{T}$ é o número de folhas e $\widehat y_{R_k}$ é a média das respostas na
folha $R_k$.
Se a forma lembra a regularização, não é coincidência.
\begin{itens}
  \item \pts{0,8} O que se obtém nos dois extremos, $\alpha=0$ e
        $\alpha\to\infty$? Diga em cada caso quantas folhas tem a subárvore
        vencedora.
  \item \pts{1,3} Considere uma folha com RSS igual a $S$, que pode ser dividida
        em duas com RSS $S_a$ e $S_b$. Escreva a condição sobre
        $S-(S_a+S_b)$ para que a divisão valha a pena segundo $C_\alpha$, e
        interprete: em que unidade o $\alpha$ está cobrando?
  \item \pts{0,8} Uma árvore tem RSS total $120$ com $6$ folhas; podando-a para
        $4$ folhas, o RSS sobe para $140$. Para quais valores de $\alpha$ a podada
        é preferível?
  \item \pts{1,1} Por que a poda usa um critério diferente do crescimento, em vez
        de simplesmente parar de crescer quando o ganho fica pequeno? Dê um
        exemplo de estrutura nos dados que um critério de parada perderia e a poda
        recupera.
\end{itens}

\begin{solucao}
Um fato prepara o exercício todo: \textbf{dividir uma região nunca aumenta o RSS}. Se
a região $R$ se divide em $R_a$ e $R_b$, então
\[
\begin{aligned}
  S=\sum_{i\in R}(Y_i-\bar Y_R)^2
   &=\sum_{i\in R_a}(Y_i-\bar Y_R)^2+\sum_{i\in R_b}(Y_i-\bar Y_R)^2\\
   &\ge\sum_{i\in R_a}(Y_i-\bar Y_{R_a})^2+\sum_{i\in R_b}(Y_i-\bar Y_{R_b})^2
   =S_a+S_b ,
\end{aligned}
\]
porque, dentro de cada parte, a média é a constante que minimiza a soma de quadrados
--- é o melhor preditor constante da Lista Teórica~01, na versão amostral.

\smallskip
\textbf{(a)} \emph{Com $\alpha=0$}, $C_0(T)$ é só o RSS. Toda subárvore $T\subseteq T_0$
se obtém de $T_0$ juntando folhas, e pelo fato acima isso nunca diminui o RSS; logo a
árvore inteira $T_0$ minimiza $C_0$, com todas as suas $\abs{T_0}$ folhas --- no limite
do crescimento, uma observação por folha e RSS zero. (Uma subárvore com o mesmo RSS de
$T_0$ empataria com ela.)

\emph{Quando $\alpha\to\infty$}, seja $S_{\mathrm{raiz}}$ o RSS da raiz sozinha, isto é,
a soma de quadrados total. A raiz custa $S_{\mathrm{raiz}}+\alpha$, e qualquer árvore com
$\abs{T}\ge2$ folhas custa pelo menos $2\alpha$. Assim que $\alpha>S_{\mathrm{raiz}}$, a
raiz vence: \textbf{uma folha só}, que prediz a média de todas as respostas. Não é
preciso $\alpha$ infinito; basta passar da soma de quadrados total.

Entre os dois extremos, as subárvores vencedoras formam uma sequência finita e
encaixada, de $T_0$ até a raiz, podando um ramo de cada vez --- e é essa sequência que a
validação cruzada percorre para escolher o $\alpha$.

\smallskip
\textbf{(b)} Dividir a folha troca $1$ folha por $2$: o RSS cai de $S$ para $S_a+S_b$, e
o custo de complexidade sobe $\alpha$. A divisão reduz $C_\alpha$ se, e só se,
\[
  \big(S_a+S_b\big)+\alpha\big(\abs{T}+1\big)<S+\alpha\abs{T}
  \iff
  \boxed{\;S-\big(S_a+S_b\big)>\alpha\;}
\]
(na igualdade, é indiferente). \textbf{A divisão precisa reduzir o RSS em mais de
$\alpha$}. O $\alpha$ é literalmente o \emph{preço de uma folha}, cobrado na mesma
unidade do RSS: se $Y$ está em reais, $\alpha$ está em reais ao quadrado --- o que
explica, de passagem, por que um valor de $\alpha$ não se transporta de um problema
para outro de escala diferente.

\smallskip
\textbf{(c)} A árvore grande custa $120+6\alpha$; a podada, $140+4\alpha$. A podada é
preferível quando
\[
  140+4\alpha<120+6\alpha
  \iff 20<2\alpha
  \iff \boxed{\;\alpha>10\;},
\]
com empate em $\alpha=10$ e as seis folhas mantidas abaixo disso.

O $10$ é $20/2$: o aumento de RSS \emph{por folha removida}. E a conta é sempre essa.
Se uma árvore tem $L$ folhas a menos que outra e RSS maior em $\Delta$, a menor vence
exatamente quando
\[
  \alpha>\frac{\Delta}{L}.
\]
O item (b) é o caso $L=1$. A regra dá um jeito rápido de estimar de cabeça qual $\alpha$
separa duas subárvores --- e, mais importante, mostra que a poda julga \emph{grupos} de
cortes pelo que eles rendem juntos, por folha.

\smallskip
\textbf{(d)} Porque o critério de parada é \textbf{míope}: julga cada corte pelo que ele
rende sozinho, e um corte que rende pouco pode ser indispensável por \emph{viabilizar}
os seguintes. A poda julga ramos inteiros --- é o critério de (c) --- e por isso enxerga
o que só aparece depois de dois ou mais cortes.

O exemplo clássico é a \textbf{interação}. Tome $X_1$ e $X_2$ independentes e uniformes
em $[-1,1]$, e $Y=\operatorname{sinal}(X_1X_2)$, sem ruído, com $n$ observações. A média
de $Y$ é zero, e $S_{\mathrm{raiz}}=\sum_iY_i^2=n$.
\begin{itemize}[leftmargin=1.2cm]
  \item \emph{Um corte sozinho não rende nada.} Depois de um corte $X_1\le t$, os dois
        lados continuam com média zero, porque, pela independência,
        $\E[\operatorname{sinal}(X_1X_2)\mid X_1\in A]
        =\E[\operatorname{sinal}(X_1)\mid X_1\in A]\cdot\E[\operatorname{sinal}(X_2)]=0$;
        o mesmo vale para cortes em $X_2$. Na população, o ganho de qualquer primeiro
        corte é exatamente zero; na amostra, é só flutuação. Um critério de parada com
        qualquer limiar positivo devolve a raiz sozinha.
  \item \emph{Dois níveis de cortes resolvem tudo.} Cortando $X_1$ em $0$ e depois
        $X_2$ em $0$ dos dois lados, as quatro folhas são os quadrantes, cada uma com
        $Y$ constante: RSS zero. Pelo critério de (c), essa subárvore de $4$ folhas vence
        a raiz sempre que $\alpha<n/3$ (ganho $n$ por três folhas a mais) e a de $2$
        folhas sempre que $\alpha<n/2$.
\end{itemize}
Quem cresce primeiro e poda depois chega aos quadrantes e os mantém para qualquer
$\alpha$ razoável; quem para cedo nunca chega lá. (Se o primeiro corte da árvore
gananciosa cair fora do zero --- todos empatam, afinal ---, ela gasta mais folhas para
separar os quadrantes, mas a conclusão não muda: o ramo inteiro rende uma fração fixa de
$n$ por folha, e o primeiro corte, sozinho, zero.)

É por isso que crescer e podar usam critérios diferentes de propósito, e por que não se
deve confundir o $\alpha$ da poda com um limite de profundidade: limitar a profundidade
é, ele próprio, um critério de parada.
\end{solucao}
